\section{The smallest flat example}
\label{sec:dual-number-obstruction}

The general construction has a rank-two example that can be checked
without a companion matrix.  Put $O=\Z_{(2)}$ and
$A=O[\epsilon]/(\epsilon^2)$.

For $p=2$, the cyclotomic matrix is $(-1)$.  The smallest lattice is
\[
 L=A^2,\qquad
 T=\begin{pmatrix}1&\epsilon\\\epsilon&-1\end{pmatrix}.
\]
Its reduced even Tate class is the trivial first coordinate and its odd
Tate class is the sign second coordinate.  The two extension boundaries
send these basis classes to one another, so $u=I$ and
$\nu_{t+1}(L)=1$.  Every scalar-extended integral lattice has invariant zero.
In particular
\[
 [L]\notin\mathbb Z[A]+\mathbb Z[A[C_2]],\qquad
 0\ne[L]-[A[C_2]],\qquad
 2([L]-[A[C_2]])=0.
\]
Since $S_2=C_2$, the degree-two Young permutation classes are precisely
$\rho_2=[A]$ and $\rho_1^2=[A[S_2]]$.  Consequently
\[
 \mathbb Z[\rho_1,\rho_2,\ldots]\longrightarrow R(A)
\]
is not surjective.

The ring $A$ is noetherian, local, of characteristic zero, free of rank two
over $\mathbb Z_{(2)}$, and has regular reduction $A/(\epsilon)=O$.  Its
only square roots of one are $\pm1$: writing a unit as $a+b\epsilon$ gives
$a^2=1$ and $2ab=0$, forcing $a=\pm1$ and $b=0$.  Moreover
$A[1/2]=\mathbb Q[\epsilon]/(\epsilon^2)$ is local, so its Picard group is
zero.  The rank-two example becomes trivial plus sign after inverting two;
explicit eigenvectors are $(1,\epsilon/2)$ and $(-\epsilon/2,1)$.
Thus neither extra roots of unity nor a nontrivial generic line bundle is
responsible for this obstruction.  It is detected by the square of the
boundary attached to the nilpotent extension.

\subsection{A relative exact sequence for involutions}

The following elementary sequence will also be useful over ramified bases.
Let $P=A[C_2]$.  For any $A[C_2]$-lattice $M$ of
$A$-rank $d$, with involution $T$, there is an admissible sequence
\begin{equation}
\label{eq:dual-relative-resolution}
 0\longrightarrow M\otimes\mathrm{sgn}
 \longrightarrow P\otimes_A M_{\mathrm{underlying}}
 \longrightarrow M\longrightarrow0.
\end{equation}
Here the group acts only on the first factor of the middle term, which is
$P^d$.  In the two copies of $M$ corresponding to $1,s$, the maps are
\[
 v\longmapsto(-Tv,v),\qquad
 (u,v)\longmapsto u+Tv.
\]
The latter has the $A$-linear section $m\mapsto(m,0)$; the group action on
the kernel parameter $v$ is $-T$, as required.  
